\clearpage
\appendix
\renewcommand*{\sectionformat}{\textcolor{cgteal}{Appendix \thesection}\hspace{0.7em}}
\addtocontents{toc}{\protect\vspace{6pt}}
\addtocontents{toc}{\protect\noindent{\protect\sffamily\protect\small\protect\color{cgmuted}Appendices, outside the page count of the report}\protect\par}

\section{Requirements of the brief and where each one is met}
\label{app:trace}
\vspace{-0.45\baselineskip}{\color{cgrule}\rule{\textwidth}{0.6pt}}\par\vspace{0.15\baselineskip}
The table follows the brief of Projet Applicatif 2 line by line. The same checklist, with direct links to the
files, is \repofile{docs/18_requirements_checklist.md}.

{\tablefont
\begin{xltabular}{\textwidth}{P{4.95cm}P{1.35cm}LP{1.3cm}}
\toprule
\headrow\hd{Requirement of the brief} & \hd{Section} & \hd{Evidence in the repository} & \hd{Status} \\
\midrule
\endfirsthead
\toprule
\headrow\hd{Requirement of the brief} & \hd{Section} & \hd{Evidence in the repository} & \hd{Status} \\
\midrule
\endhead
\bottomrule
\endfoot
\multicolumn{4}{l}{\textbf{\color{cgnavy}Common framework}} \\*
Project over two months, eight weeks & \ref{sec:project} & Week 1 from 21 September 2026; final day 10 November 2026 & Done \\
Detailed weekly progress report to the referent & \ref{sec:project} & Sent by email; progress reports of 5 and 16 October; \code{docs/14}, \code{docs/16} & Done \\
Demonstration of the realisation & \ref{sec:valo} & \code{scripts/demo.py} (15 steps), dashboard, website; \code{docs/15}, \code{docs/17} & Ready \\
Conference proposal: title, abstract, targeted venue & \ref{sec:valo} & \code{docs/13}; full paper in \code{paper/} & Done \\
\midrule
\multicolumn{4}{l}{\textbf{\color{cgnavy}Pedagogical objectives}} \\*
Understand the threats specific to AI supply chains & \ref{sec:scoping}, \ref{sec:soa} & \code{docs/02} section 2, \code{docs/04} & Done \\
Analyse the risks for model provenance, dataset lineage and pipeline integrity & \ref{sec:threat} & \code{docs/04}: assets, actors, trust boundaries, STRIDE, residual risks & Done \\
Study cryptographic signing, hash chains and Merkle trees & \ref{sec:soa}, \ref{sec:design} & \code{docs/02} sections 3 and 4, \code{docs/07} & Done \\
Design a coherent attestation architecture & \ref{sec:design} & \code{docs/05}, \code{docs/06}, \code{docs/07} & Done \\
Implement a functional system & \ref{sec:impl} & Source code, 168 tests, two CI workflows passing & Done \\
Evaluate the results objectively & \ref{sec:eval}, \ref{sec:discussion} & \code{results/}, \code{docs/08}, including the measured limit & Done \\
\midrule
\multicolumn{4}{l}{\textbf{\color{cgnavy}Phase 1: scoping and problem definition}} \\*
Scope for the attestation of AI artifacts in an enterprise; a clear cybersecurity problem & \ref{sec:intro}, \ref{sec:scoping} & \code{docs/01} sections 1, 2 and 6 & Done \\
Assets: datasets, pretrained models, fine-tuned weights, training scripts, inference APIs & \ref{sec:scoping} & \code{docs/01} section 3, \code{docs/04} section 1 & Done \\
Threats: data poisoning, backdoor injection, model substitution, pipeline compromise & \ref{sec:scoping}, \ref{sec:threat} & \code{docs/01} section 4, \code{docs/04} & Done \\
Deliverable: scoping note of 2 to 3 pages & \ref{sec:scoping} & \code{docs/01\_scoping\_note.md} & Done \\
\midrule
\multicolumn{4}{l}{\textbf{\color{cgnavy}Phase 2: state of the art and critical analysis}} \\*
Attacks: direct poisoning, backdoored base models, environment poisoning & \ref{sec:soa} & \code{docs/02} sections 2.1 to 2.3 & Done \\
Existing solutions: cryptographic signing, SLSA for ML, model cards, Sigstore & \ref{sec:soa} & \code{docs/02} sections 3 to 5 & Done \\
Compare open-source, industrial and academic approaches; identify the limitations & \ref{sec:soa} & \code{docs/02} sections 8 and 9, \code{docs/03} & Done \\
Deliverables: synthesis of 3 to 5 pages, comparative table, argued justification of the approach & \ref{sec:soa} & \code{docs/02} (justification in section 10), \code{docs/03} & Done \\
\midrule
\multicolumn{4}{l}{\textbf{\color{cgnavy}Phase 3: solution design}} \\*
Provenance and integrity architecture for AI artifacts & \ref{sec:design} & \code{docs/05} & Done \\
Signing, verification and revocation mechanisms & \ref{sec:design} & \code{docs/07}; \code{chainguard/trustlog.py} & Done \\
Justified technical choices & \ref{sec:soa}, \ref{sec:design} & \code{docs/05} section 7, \code{docs/02} section 10 & Done \\
Deliverables: architecture diagrams, functional and technical description, attack scenarios & \ref{sec:design} & \code{docs/05} (six diagrams), \code{docs/06}, \code{docs/07}, \code{docs/08} & Done \\
Indicative architecture: hash chain or Merkle tree & \ref{sec:design} & Hash-linked manifests and row-level Merkle roots for datasets & Done \\
Indicative architecture: signing of models and datasets at each pipeline stage & \ref{sec:design} & Ed25519 manifests for datasets, base model, fine-tuned model and deployment record & Done \\
Indicative architecture: blockchain anchoring (optional) & \ref{sec:soa} & Not built; a signed, hash-chained, pinned trust log instead, justified in \code{docs/02} sections 4 and 10 & Optional, not built \\
Indicative architecture: CI/CD integration for automated attestation & \ref{sec:impl} & \code{.github/workflows/attest-and-verify.yml} & Done \\
Indicative architecture: verification API for deployment gateways & \ref{sec:impl} & \code{api/} and \code{deployment/gateway.py} & Done \\
\midrule
\multicolumn{4}{l}{\textbf{\color{cgnavy}Phase 4: implementation and experimentation}} \\*
Simulated ML pipeline & \ref{sec:impl} & \code{ml\_pipeline/} & Done \\
Attestation across at least three pipeline stages & \ref{sec:design} & Four stages & Done \\
Documented code and scripts & \ref{sec:impl} & \code{README.md}, docstrings, \code{scripts/}, \code{docs/07} & Done \\
Description of the test environment & \ref{sec:impl} & \code{docs/09}, \code{results/validation\_logs/environment.txt} & Done \\
Experimentation journal & \ref{sec:eval} & \code{docs/09\_experimentation\_journal.md} & Done \\
Technologies: Python, PyTorch; Sigstore, cosign or custom signing with OpenSSL; Merkle tree library; Docker, GitHub Actions; Hyperledger Fabric (optional) & \ref{sec:impl} & Python and PyTorch; custom Ed25519 signing on OpenSSL and an optional Sigstore job; own RFC 6962 implementation; \code{docker/}; \code{.github/workflows/}; Hyperledger not used & Done \\
\midrule
\multicolumn{4}{l}{\textbf{\color{cgnavy}Phase 5: testing, evaluation and critical analysis}} \\*
Test against realistic attack scenarios & \ref{sec:eval} & 48 variants in seven experiments, \code{attacks/}, \code{docs/08} & Done \\
Measure detection rates and overhead & \ref{sec:eval} & \code{results/attacks/}, \code{results/metrics/} & Done \\
Deliverables: test results as logs, tables and graphs & \ref{sec:eval} & \code{results/validation\_logs/}, \code{results/attacks/summary.csv}, \code{results/figures/}, \code{report/figures/} & Done \\
Critical analysis of the guarantees and the limitations; improvement paths & \ref{sec:discussion} & Paper sections V and VI, \code{docs/04} (residual risks) & Done \\
The four example tests & \ref{sec:eval} & Table~\ref{tab:brieftests} of this report & Done \\
\midrule
\multicolumn{4}{l}{\textbf{\color{cgnavy}Phase 6: reporting and valorisation}} \\*
Final report of about 30 pages & all & This document; source in \code{report/} & Done \\
Oral defence & \ref{sec:valo} & Presentation guide (\code{docs/17}) and slides prepared & Ready \\
Demonstration or use case & \ref{sec:valo} & Demonstration script, dashboard, website & Done \\
Conference proposal & \ref{sec:valo} & \code{docs/13}, \code{paper/} & Done \\
\midrule
\multicolumn{4}{l}{\textbf{\color{cgnavy}Expected results}} \\*
A functional attestation system; a complete engineering approach with threat modelling; a critical analysis of the results; work that can be valorised in an enterprise security context & all & The repository, documents \code{01} to \code{09}, this report, the paper and the website & Done \\
\end{xltabular}}

\needspace{14\baselineskip}
\section{The 48 variants}
\label{app:variants}
\vspace{-0.45\baselineskip}{\color{cgrule}\rule{\textwidth}{0.6pt}}\par\vspace{0.15\baselineskip}
Reference run of 9~November 2026. For every variant the observed status equals the status predicted before the
run. The time is that of the verification of the attacked artifact with its chain.

\makeatletter{\tablefont
\begin{xltabular}{\textwidth}{P{1.3cm}LP{1.75cm}P{2.45cm}r}
\toprule
\headrow\hd{Variant} & \hd{Manipulation} & \hd{Status} & \hd{Checks that fired} & \hd{Time (ms)} \\
\endfirsthead
\toprule
\headrow\hd{Variant} & \hd{Manipulation} & \hd{Status} & \hd{Checks that fired} & \hd{Time (ms)} \\
\endhead
\bottomrule
\endfoot
\@@input sections/gen_variants.tex 
\end{xltabular}}\makeatother

\needspace{20\baselineskip}
\section{A signed manifest}
\label{app:manifest}
\vspace{-0.45\baselineskip}{\color{cgrule}\rule{\textwidth}{0.6pt}}\par\vspace{0.15\baselineskip}
The manifest of the fine-tuned model of the reference run, as stored on disk. Three provenance fields are
shortened here for space (marked by three dots), so this printed version is not the signed payload itself. The
two parents are the base model and the fine-tuning dataset, each named by its manifest digest.

\begin{lstlisting}
{
  "payload_type": "application/vnd.chainguard.manifest+json; version=1",
  "payload": {
    "schema_version": "1.0",
    "artifact_id": "9f25b977-0aa0-4aa8-a149-4f19d79f3ee2",
    "artifact_name": "digits-finetuned-cnn",
    "artifact_type": "finetuned_model",
    "version": "1.0.0",
    "pipeline_stage": "stage3_finetune",
    "subject": {
      "uri": "artifacts/models/digits-finetuned-cnn-v1.0.0.safetensors",
      "sha256": "1f1448ea86c028d9d36e80b764eca9a52368fdf9847cb901db394f338b16a0ae",
      "size_bytes": 153944,
      "media_type": "application/safetensors",
      "merkle": null
    },
    "parents": [
      { "artifact_id": "a85de012-e158-4221-85e7-826ff7749d5f",
        "manifest_digest": "sha256:1f930c8538ccf7ed0ef95813eed5d8fad5741bad7956e8fab1faf6d8ea82fd21",
        "relation": "fine_tuned_from" },
      { "artifact_id": "3c4e300b-12fe-4883-89c6-cdb3ce96758a",
        "manifest_digest": "sha256:6959cdaa0188c84cf87e0b9b28dfa6d84c0e1f38b68193f5fb49133f212048ea",
        "relation": "trained_on" }
    ],
    "provenance": {
      "source": "internal fine-tuning pipeline",
      "builder": "ml_pipeline/finetune.py",
      "parameters": { "epochs": 8, "lr": 0.0005, "batch_size": 64, "seed": 42 },
      "metrics": { ... },
      "input_verification": { ... },
      "environment": { "python": "3.13.16", "chainguard": "0.4.0", "torch": "2.15.0+cu132", ... }
    },
    "semantic_checks": [],
    "created_at": "2026-11-09T09:26:46Z",
    "signer": {
      "identity": "training_pipeline@chainguard.local",
      "role": "training_pipeline",
      "key_id": "ed25519:cec6a3e317c4a128"
    }
  },
  "signatures": [
    { "key_id": "ed25519:cec6a3e317c4a128", "algorithm": "Ed25519",
      "sig": "n5sYF5FF/diG4Vp88bYcJcsciUHMH9cHh+N8uvt+vXnMwpxp/DCfUaGCEpX3G0yhkzudUNgZMkF3Fuban7syBA=="
    }
  ],
  "manifest_digest": "sha256:133c0274c7b49c52fb8a58504bb8f090a7a28853d8f1e0be98ffd73004f1e0e9"
}
\end{lstlisting}

The last field is a convenience copy. The verifier never trusts it: it recomputes the digest from the payload.

\section{Reproducing the results}
\label{app:repro}
\vspace{-0.45\baselineskip}{\color{cgrule}\rule{\textwidth}{0.6pt}}\par\vspace{0.15\baselineskip}
From a clone of the repository, with Python 3.11 or later:

\begin{lstlisting}
make install          # dependencies, PyTorch for CPU
make keys             # one Ed25519 key pair per role
make pipeline         # the four-stage attested pipeline and the deployment gateway
make attacks          # the 48 variants of the seven experiments
make benchmark        # cost of hashing, signing and verification
make test             # the 168 tests
python -m scripts.e2_finetune_poison      # supplementary run: poison in the fine-tuning set
python -m scripts.seed_robustness         # supplementary run: ten other seeds
python report/figures/make_figures.py     # the charts of this report
make demo             # the ten-minute demonstration
make standalone       # the dashboard, on port 8501
docker compose -f docker/docker-compose.yml up --build    # pipeline, API and dashboard in containers
make validate LABEL="<machine>"           # the whole validation sequence, with its logs
\end{lstlisting}

\section{Deliverables and where to find them}
\label{app:deliverables}
\vspace{-0.45\baselineskip}{\color{cgrule}\rule{\textwidth}{0.6pt}}\par\vspace{0.15\baselineskip}
{\tablefont
\begin{tabularx}{\textwidth}{P{4.0cm}L}
\toprule
\headrow\hd{Deliverable} & \hd{Location} \\
\midrule
This report, PDF and LaTeX source & \repodir{report} \\
Website & \url{\site} \\
Source code, tests, Docker and CI files & \url{\repo} \\
Conference paper, PDF and LaTeX source & \repodir{paper} \\
Conference proposal & \repofile{docs/13_conference_proposal.md} \\
Scoping note, state of the art, comparative table, threat model, architecture, functional and technical descriptions, attack scenarios, experimentation journal & \repodir{docs}, documents \code{01} to \code{09} \\
Requirements checklist & \repofile{docs/18_requirements_checklist.md} \\
Results: logs, tables, figures, archived earlier runs & \repodir{results} \\
Demonstration runbook & \repofile{docs/15_professor_demo_10min.md} \\
Earlier progress reports of 5 and 16 October 2026 & Folder \emph{Previous reports} of the hand-in package \\
\bottomrule
\end{tabularx}}
